\documentclass{article}
\usepackage{graphicx} % Required for inserting images
\usepackage{tcolorbox}
\usepackage{amsthm}
\usepackage{amsmath}
\usepackage{amssymb}
\usepackage{enumitem}
\usepackage{hyperref}
\usepackage{braket}

\title{Quantum Computation Notes \\ \large Regev's Algorithm}
\author{}
\date{}
\begin{document}
\maketitle
We begin with some useful definitions and theorems.\\
\textbf{Definition:} Given $n$ linearly independent vectors (referred to as the basis), $b_1, \dots, b_n \in R^{m}$ the lattice generated by them is defined as $L(b_1, \dots, b_n) = \{ \sum x_i b_i | x_i \in \mathbb{Z}\}$. If we define $B$ as the $m \times n$ matrix whose columns are $b_1, \dots, b_n$ the lattice generated by $B$ is  $L(B) = L(b_1, \dots, b_n) = \{ Bx | x \in \mathbb{Z}^n \}$. $L$ is a lattice if and only if $L$ is a discrete subgroup of $(R^n, +)$. Alternatively a lattice $L = B\mathbb{Z}^n$.\\
\textbf{Definition:} For matrix $B$, $P(B) = \{ Bx | x \in [0,1)^n\}$ is the fundamental parallelepiped of $B$, the area spanned by the basis vectors with scalar coefficient 1.\\
\textbf{Definition:} The \textit{determinant} of the lattice is $\det(L) = \sqrt{\det(B^TB)}$, the volume of $P(B)$.\\
\textbf{Definition:} The dual of a lattice $\mathcal{L}$ is $\mathcal{L}^*$ of all vectors in $x \in \text{span }\mathcal{L}$ such that $\braket{x,y}$ is an integer for all $y \in \mathcal{L}$.\\

\textbf{Theorem 1:} Suppose $N$ is an $L$ bit composite number and $x$ is a non-trivial solution to $x^2 \equiv 1 \pmod N$ in the range $1\leq x \leq N$, with nontrivial meaning that $x \neq 1$ and $x \neq N-1$. Then at least one of gcd$(x\pm1, N)$ is a non-trivial factor of $N$ that can be computed in $O(L^3)$ operations.\\
\textbf{The result:} Fix $n-$bit number $N \leq 2^n$ to be factorized. For some $d>0$, let $b_1, \dots, b_d$ be some small $O(\log d)$-bit integers coprime to $N$ (say $b_i$ is the $i$th prime) and let $a_i = b_i^2$. Define lattice:
$$\mathcal{L} = \{ (z_1, \dots, z_d) \in \mathbb{Z}^d | (\prod_i b_i^{z_i})^2 \equiv 1 \bmod N\}$$
$$=\{(z_1, \dots, z_d) \in \mathbb{Z}^d | \prod_i a_i^{z_i} \equiv 1 \bmod N\}$$
Define sublattice $\mathcal{L}_0$:
$$\mathcal{L}_0 = \{ (z_1, \dots, z_d) \in \mathbb{Z}^d | \prod_i b_i^{z_i} \in \{-1, 1\} \bmod N \} \subseteq L$$
Given avector $z \in \mathcal{L} \backslash \mathcal{L}_0$ we have that $b = \prod_i b^{z_i} \bmod N$ is a square root of unity modulo $N$ yet it is a non-trivial square root of 1 (not equal to $\pm 1 \bmod N$, by $z \not \in \mathcal{L}_0$) (theorem 1).
In this case, $N$ divides the product $(b-1)(b+1)$ but does not divide either of the terms and we must have that gcd$(b-1, N)$ is a non-trivial factor of $N$ as desired. \\Therefore the objective of the algorithm is to find a vector in $\mathcal{L} \backslash \mathcal{L}_0$. \\ 
\textbf{Result (Algorithm):} There is a classical polynomial time algorithm that does the following: \\
As input, it receives an $n-$bit number $N$ and $O(\log n)$-bit numbers $b_1, \dots, b_d$ where $d= \sqrt{n}$. The algorithm performs $\sqrt{n} + 4$ calls to a quantum circuit of size $O(n^{3/2} \log n)$ that uses $O(n^{3/2})$ qubits.
This relies on a conjecture for the choice of $b_i$, which is that:\\
\textbf{Conjecture 1.2:} Let $N$ be an odd $n-$bit number that is not a prime power. Then, with probability at least half, for $d = \sqrt{n}$, if $b_1, \dots, b_d$ are randomly chosen $O(\log n)$ bit prime numbers, $\mathcal{L}\backslash \mathcal{L}_0$ contains a vector of norm at most  $T = e^{O(\sqrt{n})}$.

\textbf{High Level Overview of Algorithm:}\\
The core of the algorithm is the next section (the quantum procedure).
\begin{enumerate}
	\item Initialize two registers, in discrete Gaussian states of some radius $R$.
	\item Create superposition over $\ket{z}, z \in \mathbb{Z}^d$.
	\item Compute $\prod_i a_{i}^{z_i} \bmod N$ in a new register (on small numbers, $O(n^{1.5})$, so we only do $n-$bit multiplication are squaring and multiplication operations. We can measure that register now, fixing the first register into a superposition of a random coset of $L$.\\
	\item Apply QFT and measure, obtaining vectors in $\mathcal{L}^*$.
	\item Use continued fractions to recover $\mathcal{L}$.
\end{enumerate}
\section*{The Quantum Procedure:}\\
Fix $d > 0, a_1, \dots, a_d$ and $\mathcal{L}$. Let $R> \sqrt{2d}$ be another parameter chosen later and let $D$ be a power of two larger than $R$, say $D \in [ 2\sqrt{d} R, 4\sqrt{d} R)$.\\
The quantum procedure outputs a uniform sample from $L^* \backslash \mathbb{Z}^d$ (a set of size $\det \mathcal{L}$) perturbed by Gaussian noise of $1/R$ and discretized to grid $\{0,1/D, \dots, (D-1)/D \}^d$. (what exactly is a discrete Gaussian state? what is the significance of its radius?)\\ 
In more detail about the noise: the output will be witihn statistical distance $1/poly(d)$ of the distribution $Q$ supported on the discretized grid and whose mass at point $w$ is:
$$Q(w) \equiv (\det \mathcal{L})^{-1} \sum_{v \in \mathcal{L}^* / \mathcal{Z}^d} Q_v(w)$$
where for $$v \in \mathcal{L}^* / \mathbb{Z}^d, Q_v(w) \equiv \frac{\rho_{1/\sqrt{2R} (v+ w + Z^d)}}{\rho_{1/\sqrt{r}} (v- D^{-1} \mathbb{Z}^d)}$$
for Gaussian function $\rho_s(z) = exp(-\pi ||z||^2 / s^2)$.\\
\begin{enumerate}
	\item Approximate within $1/poly(d)$ the state proportional to:
		$$\sum_{z \in \{ -D/2, \dots, D/2-1\}} \rho_R(z) \ket{z}$$
		essentially a superposition of states, analogous to Hadamard step except its a superposition of Gaussian states instead. After the first $d$ qubits, just set to $\ket{+} = \ket{0}+\ket{1}/\sqrt{2}$. Uses quantum circuit of size $d(\log D + \text{poly}(\log d))$.
	\item Compute $h(z) := \prod_i a_i^{z_i + D/2} \bmod N$ in a new register. Then, the new state is:
		$$\sum_{z \in \{ -D/2, \dots, D/2-1\}}^d \rho_r(z) \ket{z} \ket{h(z)}$$. This is calculated in a way similar to a binary tree, like a repeated squaring building a tree. This can be done in $O(\log D (d \log^3 d + n \log n)$ steps.
\item Apply $QFT$ to $\ket{z}$ register, and then output vector in $\{ 0, 1/D, \dots, (D-1)/D\}^d$ obtained by measuring register and dividing by $D$. This is similar to switching to Fourier basis step in usual. To elaborate, consider $h'(z) := \prod_i a_i^{z_i} \bmod N \iff h(z) = \prod_i a_i^{D/2} h'(z) \bmod N$ and thus by applying some change we can obtain:
	$$\sum_{z \in \{ -D/2, \dots, D/2-1\}}^d \rho_r(z) \ket{z} \ket{h(z)} = \sum_{z\in \{ -D/2, \dots, D/2-1\}}^d \rho_R(z) \ket{z} \ket{h'(z)}$$
	Since $h'$ is homomorphism from $Z^d$ to $G= (Z_n, \times)$ (mult group of int. mod N) with kernel $L$ since if its on the lattice then the remainder will be zero, by isomorphism theorem ($G/ker G \cong img(G)$) we can change the second reigster to some element in $Z^d/L$, so the input to the final step of QFT and measure is;
	$$\sum_{e \in Z^d/L}  \sum_{z \in (L+e) \cap [-D/2, D/2)^d} \rho_r(z) \ket{z} \ket{e}$$
	Taking $\epsilon = 1/poly(d)$ we get circuit size $O(d \log D (\log \log D + \log d))$ thus the quantum procedure uses circuit of size:
	$$O(d \log^3 D  + d \log^2 D + n \log n) + d poly(\log d))$$
	and otuputs a point $w \in [0,1)^d$ that is within distance $\sqrt{d}/\sqrt{2}R$ of a uniformly chosen $v \in L^*/Z^d$. With our choice of parameters, this becomes $O(n \log n \log D)$.
\end{enumerate}
\section*{Recovering a Lattice from Noisy Samples from the Dual}
dog
\end{document}
